\documentclass[11pt,a4paper]{amsart}

\usepackage{amsmath,amssymb,amsthm,booktabs,microtype,hyperref}
\usepackage[margin=1in]{geometry}
\sloppy

\newtheorem{theorem}{Theorem}[section]
\newtheorem{proposition}[theorem]{Proposition}
\newtheorem{lemma}[theorem]{Lemma}
\newtheorem{corollary}[theorem]{Corollary}
\theoremstyle{definition}
\newtheorem{definition}[theorem]{Definition}
\newtheorem{remark}[theorem]{Remark}

\newcommand{\R}{\mathbb{R}}
\newcommand{\C}{\mathbb{C}}
\newcommand{\Q}{\mathbb{Q}}
\newcommand{\Z}{\mathbb{Z}}
\newcommand{\RP}{\mathbb{RP}}
\DeclareMathOperator{\Hess}{Hess}
\DeclareMathOperator{\tr}{tr}
\DeclareMathOperator{\sgn}{sgn}
\DeclareMathOperator{\diag}{diag}

\title[Ninety-Four Positive Electrostatic Equilibria for Eleven Charges]{%
  Ninety-Four Nondegenerate Equilibria for Eleven Positive Point Charges:\\
  Breaking the Spherical 5-Design Barrier via Multi-Shell Kelvin Designs,\\
  Recursive Harmonic Nesting, and Pandrosion V31 / \texttt{pandsysroots} Certification}

\author{Ivan Besevic}
\address{ESIEE-IT, France}
\email{ivan.besevic@edu.esiee-it.fr}

\subjclass[2020]{Primary 31B05, 58K05; Secondary 14P05, 65G20, 68V20}
\keywords{Maxwell's conjecture, electrostatic equilibria, Coulomb potential,
Kelvin design, harmonic polynomials, Morse index, Krawczyk operator,
Pandrosion projective radical enclosure, \texttt{pandsysroots}}

\begin{document}

\begin{abstract}
Let $M^+(n)$ denote the supremum number of finite, nondegenerate critical points
of the three-dimensional Coulomb potential
$V(x) = \sum_{j=1}^n q_j / \|x - a_j\|$ generated by $n$ distinct strictly
positive point charges $q_j > 0$ in $\R^3$. Following the July 2026 refutation
of Maxwell's $(n-1)^2$ conjecture by Arathoon, Ball, and Kvalheim
($M^+(5) \ge 24$, together with a $+20$ axial-pair lift $M^+(n+2) \ge M^+(n)+20$)
and Kintali's subsequent positive counterexamples for $n \in \{4,6,8,9,10\}$
(notably $M^+(9) \ge 68$ and $M^+(10) \ge 91$), the strongest previously known
lower bound for eleven positive charges was $M^+(11) \ge 68 + 20 = 88$,
lifting to $M^+(2k+1) \ge 20k - 12$ for all $k \ge 4$.

In this paper we prove unconditionally that
\[
  M^+(11) \ge 94,
\]
and hence $M^+(2k+1) \ge 20k - 6$ for every integer $k \ge 5$ (giving
$M^+(13) \ge 114$, $M^+(15) \ge 134$, $M^+(17) \ge 154$, and improving every
odd bound for $n \ge 11$ by $+6$). Our construction introduces three new
ingredients:
(i) a \emph{multi-shell positive Kelvin $5$-design} of $11$ strictly positive
point charges in $\R^3$ that annihilates all $35$ solid spherical harmonics of
degrees $1$ through $5$ with full Jacobian rank $35$, circumventing the
classical Hardin--Sloane minimal size $n \ge 12$ for positive weighted spherical
$5$-designs on $S^2$ while simultaneously generating $23$ nondegenerate outer
equilibria ($c_1^{\mathrm{out}} = 7$, $c_2^{\mathrm{out}} = 16$);
(ii) a \emph{four-scale recursive harmonic nesting} anchored by a
$100\%$-B\'ezout $C_{3v}$ quartic harmonic polynomial $P_{27}$ (all $27/27$
critical points real and simple via an off-mirror nodal plane), extended
through degree $5$ ($+25$ equilibria, $52/64$ real) and aligned in $SO(3)$
against the $11$-charge sextic harmonic $H_6$ ($+19$ equilibria, $71/125$ real,
certified complete $125/125$ by \texttt{pandsysroots~0.34.0}); and
(iii) an explicit exact-rational configuration in $(\Z / 10^{90})^4$ whose
$94$ distinct nondegenerate critical points ($c_1 = 42$, $c_2 = 52$, matching
the exact Poincar\'e--Hopf balance $c_2 - c_1 = 10$) are certified directly
on the actual 3D Coulomb potential $V(x)$ in $41.5\,\mathrm{ms}$ by a
self-contained pure-Rust verifier using $512$-bit outward-rounded dyadic
intervals, Pandrosion~V31 two-sided order-$4$ projective radical enclosures
(zero floating-point or \texttt{libm} trust), and the 3D interval Krawczyk
operator ($\max \|I - Y_k \Hess V(B_k)\|_\infty \le 1.386 \times 10^{-21}$).
\end{abstract}

\maketitle

\footnote{Prepared with the assistance of AI tools for symbolic exploration,
code generation, and manuscript drafting; all mathematical results are
independently verified by exact symbolic identities and the accompanying
self-contained Rust certification suite.}

\section{Introduction}\label{sec:intro}

In his 1873 \emph{Treatise on Electricity and Magnetism}, James Clerk Maxwell
conjectured that the electrostatic potential
\begin{equation}\label{eq:coulomb}
  V(x) = \sum_{j=1}^n \frac{q_j}{\|x - a_j\|}, \qquad x \in \R^3 \setminus \{a_1,\dots,a_n\},
\end{equation}
of $n$ distinct point charges $a_1,\dots,a_n \in \R^3$ with nonzero strengths
$q_1,\dots,q_n \in \R \setminus \{0\}$ has at most $(n-1)^2$ isolated
equilibrium points (zeros of $\nabla V$). Formalized in modern Morse-theoretic
terms by Gabrielov, Novikov, and Shapiro~\cite{GNS2007}, let $M(n)$ denote the
supremum number of finite nondegenerate critical points of $V$ over all
$n$-charge configurations, and let $M^+(n)$ denote the corresponding supremum
restricted to strictly positive charges $q_1,\dots,q_n > 0$. Because $V$ is
harmonic on $\R^3 \setminus \{a_1,\dots,a_n\}$, every nondegenerate critical
point $x_*$ satisfies $\tr \Hess V(x_*) = \Delta V(x_*) = 0$ and therefore has
Morse index $1$ ($\det \Hess V(x_*) < 0$, one negative and two positive
eigenvalues) or Morse index $2$ ($\det \Hess V(x_*) > 0$, two negative and one
positive eigenvalue). Writing $c_1$ and $c_2$ for the respective counts, the
Poincar\'e--Hopf index theorem on $\R^3 \setminus \{a_1,\dots,a_n\}$ forces
\begin{equation}\label{eq:ph}
  c_2 - c_1 = n - 1
\end{equation}
whenever all $n$ charges are positive and all critical points are nondegenerate.

For more than a century Maxwell's bound $M(n) \le (n-1)^2$ remained open beyond
$n = 3$ (where $M(3) = 4$ is exact). In July 2026, Arathoon, Ball, and
Kvalheim~\cite{ABK2026} disproved Maxwell's conjecture even for strictly
positive charges by proving $M^+(5) \ge 24 > (5-1)^2 = 16$ via a symmetric
bifurcation of three equatorial unit charges and two tiny axial charges,
together with a local $+20$ axial-pair lift theorem:
\begin{equation}\label{eq:abk-lift}
  M^+(n+2) \ge M^+(n) + 20 \qquad (n \ge 3).
\end{equation}
Immediately afterward, Kintali~\cite{Kintali2026pos,Kintali2026quad,Kintali2026sup}
established positive counterexamples for $n \in \{4,6,8,9,10\}$:
\[
  M^+(4) \ge 11,\quad M^+(6) \ge 31,\quad M^+(8) \ge 51,\quad
  M^+(9) \ge 68,\quad M^+(10) \ge 91,
\]
using trigonal-pyramid cubic bifurcations, separated local scales, and positive
Kelvin designs realizing quintic and sextic harmonic limits, along with
asymptotic lower bounds $M^+(n) = \Omega(n^{3/2})$ for positive charges and
$M(n) \ge n^2/4 - n + 1$ for signed charges. Combining Kintali's $M^+(9) \ge 68$
with the Arathoon--Ball--Kvalheim lift~\eqref{eq:abk-lift} gave the best
previously known odd-charge lower bounds
\[
  M^+(11) \ge 88, \quad M^+(13) \ge 108, \quad M^+(15) \ge 128, \quad
  M^+(2k+1) \ge 20k - 12 \quad (k \ge 4).
\]

\subsection*{Main results}
The purpose of this paper is to break the $n = 11$ barrier directly at the
sextic harmonic level---without paying the spherical $5$-design cost of $n = 12$
charges---and certify every critical point directly on an explicit rational
$11$-charge potential with zero floating-point trust.

\begin{theorem}[Main Theorem]\label{thm:main}
There exists an explicit configuration of $11$ distinct strictly positive point
charges $(a_j, q_j)_{j=1}^{11} \subset (\Z / 10^{90})^3 \times (\Z_{>0} / 10^{90})$
whose Coulomb potential~\eqref{eq:coulomb} possesses at least $94$ distinct
finite nondegenerate critical points, consisting of $c_1 = 42$ index-$1$ saddles
and $c_2 = 52$ index-$2$ saddles ($c_2 - c_1 = 10$). In particular,
\[
  M^+(11) \ge 94.
\]
\end{theorem}

\begin{corollary}[New Odd-Charge Lower Bounds for All $n \ge 11$]\label{cor:odd}
Applying the Arathoon--Ball--Kvalheim $+20$ axial-pair lift~\eqref{eq:abk-lift}
recursively to Theorem~\ref{thm:main} yields
\[
  M^+(2k+1) \ge 20k - 6 \qquad \text{for every integer } k \ge 5.
\]
In particular, $M^+(11) \ge 94$, $M^+(13) \ge 114$, $M^+(15) \ge 134$, and
$M^+(17) \ge 154$, improving every previously known odd lower bound for
$n \ge 11$ by $+6$.
\end{corollary}

In the course of constructing Theorem~\ref{thm:main}, we also isolate two
structural phenomena of independent interest in real algebraic geometry and
potential theory:
\begin{enumerate}
  \item \emph{A $100\%$-B\'ezout degree-$4$ harmonic polynomial $P_{27}$:}
  Proposition~\ref{prop:p27} exhibits an explicit $C_{3v}$-symmetric harmonic
  polynomial of degree $4$ all $27 = 3^3$ of whose B\'ezout critical points are
  real and nondegenerate ($c_1 = 16$, $c_2 = 11$, $c_2 - c_1 = -5$), certified
  complete by \texttt{pandsysroots~0.34.0}~\cite{pandsysroots}.
  \item \emph{Circumventing the spherical $5$-design bound $n \ge 12$:}
  Whereas any positive weighted spherical $5$-design on $S^2$ requires at least
  $n \ge 12$ nodes (achieved by the regular icosahedron), allowing multiple
  radii in $\R^3$ relaxes the radial moments $\sum w_j \|b_j\|^{2m} p(b_j)$ and
  permits a full-rank, well-separated $11$-point positive Kelvin $5$-design
  that annihilates all $35$ solid harmonics of degrees $1$ through $5$ while
  simultaneously creating $23$ nondegenerate outer equilibria.
\end{enumerate}

\section{Multi-Shell Positive Kelvin $5$-Designs at $n = 11$}\label{sec:kelvin}

Let $(a_j, q_j)_{j=1}^n \subset (\R^3 \setminus \{0\}) \times \R_{>0}$ be a
configuration of positive charges. Under the Kelvin inversion
\begin{equation}\label{eq:kelvin-map}
  b_j = \frac{a_j}{\|a_j\|^2}, \qquad w_j = \frac{q_j}{\|a_j\|} > 0,
\end{equation}
which is a smooth bijection on $(\R^3 \setminus \{0\}) \times \R_{>0}$ with
inverse $a_j = b_j / \|b_j\|^2$, $q_j = w_j / \|b_j\| > 0$, the multipole
expansion of~\eqref{eq:coulomb} inside the open ball
$\|x\| < r_{\min} := \min_j \|a_j\| = (\max_j \|b_j\|)^{-1}$ takes the pure
polynomial form
\begin{equation}\label{eq:taylor}
  V(x) = \sum_{\ell=0}^{\infty} H_\ell(x), \qquad
  H_\ell(x) = \sum_{j=1}^n w_j \|b_j\|^\ell \|x\|^\ell P_\ell(\hat{b}_j \cdot \hat{x}),
\end{equation}
where $P_\ell$ is the Legendre polynomial of degree $\ell$ and
$\hat{u} = u / \|u\|$. Fixing an orthogonal basis $\{Y_{\ell,m}\}_{m}$ of
harmonic homogeneous polynomials of degree $\ell$ (solid spherical harmonics),
the addition theorem gives
\begin{equation}\label{eq:addition}
  H_\ell(x) = \sum_{m} \alpha_{\ell,m}
  \left(\sum_{j=1}^n w_j Y_{\ell,m}(b_j)\right) Y_{\ell,m}(x),
  \qquad
  \alpha_{\ell,m} = \frac{4\pi / (2\ell + 1)}{\int_{S^2} Y_{\ell,m}^2\,d\Omega}.
\end{equation}
A \emph{positive Kelvin $t$-design} of size $n$ is a point-weight set
$\{(b_j, w_j)\}_{j=1}^n$ with $w_j > 0$ such that $H_1 = \dots = H_t \equiv 0$,
a system of $\sum_{\ell=1}^t (2\ell + 1) = (t+1)^2 - 1$ polynomial equations in
$4n$ real unknowns (or $4n - 2$ essential unknowns modulo overall radial scaling
and charge scaling).

For $t = 5$, there are $35$ equations ($3+5+7+9+11 = 35$). When all nodes are
constrained to the unit sphere $\|b_j\| = 1$, classical linear-programming and
cubature bounds (Delsarte--Goethals--Seidel; Hardin--Sloane) force $n \ge 12$,
with the regular icosahedron realizing $n = 12$ and possessing $0$ outer
equilibria (all topological index $+11$ is concentrated at the origin).

When the radii $\|b_j\|$ are freed in $\R^3$, the radial constraint
$\sum w_j \|b_j\|^{2m} p(b_j) = \sum w_j p(b_j)$ disappears. A hierarchical
multi-shell layout---$7$ outer charges at $\|a_j\| \in [0.60, 1.56]$
($\|b_j\| \approx 1$), $3$ intermediate charges at $\|a_j\| \in [0.13, 0.19]$
($\|b_j\| \approx 5.5\text{--}7.6$), and $1$ inner charge at
$\|a_1\| \approx 0.0339$ ($\|b_1\| \approx 29.5$) with tiny positive weight
$q_1 \approx 4.67 \times 10^{-8}$---supplies a $w_1 \|b_1\|^5 = O(1)$ degree-$5$
correction with negligible lower-degree weight, breaking the $n \ge 12$ barrier.

\begin{proposition}[Full-Rank Multi-Shell $n = 11$ Positive Kelvin $5$-Design]\label{prop:design}
There exists a strictly positive $11$-charge configuration
$(a_j^{(0)}, q_j^{(0)})_{j=1}^{11}$ (refined to $110$-digit precision,
residual $< 6 \times 10^{-98}$) satisfying:
\begin{enumerate}
  \item $H_1 = H_2 = H_3 = H_4 = H_5 = 0$ identically at the origin;
  \item the $35 \times 44$ Jacobian of the $35$ moment equations with respect
  to $(b_j, \log w_j)$ has full rank $35$ (leaving a $7$-dimensional smooth
  nullspace modulo the two scaling symmetries);
  \item the leading harmonic term $H_6^{(0)}$ is nondegenerate with topological
  degree $\deg(\nabla H_6^{(0)} / \|\nabla H_6^{(0)}\| : S^2 \to S^2) = +1$;
  \item the unperturbed potential $V^{(0)}$ already admits $23$ distinct
  nondegenerate outer equilibria in the annulus $0.036 \le \|x\| \le 0.801$,
  with Morse counts $c_1^{\mathrm{out}} = 7$, $c_2^{\mathrm{out}} = 16$
  ($c_2^{\mathrm{out}} - c_1^{\mathrm{out}} = +9 = 10 - \deg(\nabla H_6^{(0)})$).
\end{enumerate}
\end{proposition}

\begin{remark}[Exact Harmonicity of $Y_{5,3c}$]
In polynomial coordinates $(x,y,z)$ with $\rho^2 = x^2 + y^2$, the degree-$5$
tesseral harmonic of order $m = 3$ arises from homogenizing
$P_5^{(3)}(u) \propto (1 - u^2)^{3/2}(9u^2 - 1)$ via $R^2 = \rho^2 + z^2$, which
yields $9z^2 - R^2 = 8z^2 - \rho^2$, so
\[
  Y_{5,3c}(x,y,z) = (x^3 - 3xy^2)(8z^2 - x^2 - y^2)
  = 8x^3 z^2 - 24 x y^2 z^2 - x^5 + 2 x^3 y^2 + 3 x y^4.
\]
All $48$ basis harmonics of degrees $1$ through $6$ used in our construction
and in the Rust verifier are verified to satisfy $\Delta Y_{\ell,m} \equiv 0$
identically in $\Z[x,y,z]$.
\end{remark}

\section{Recursive Multi-Scale Harmonic Nesting and $P_{27} \to P_{52} \to P_{71}$}\label{sec:nesting}

Given a fixed nondegenerate top-degree harmonic polynomial $H_6$ arising from
Proposition~\ref{prop:design}, the full-rank Implicit Function Theorem allows us
to prescribe \emph{arbitrary} sufficiently small lower-degree harmonic
perturbations $H_1, \dots, H_5$ by an $O(\delta)$ perturbation of the $11$
positive charges, without losing any of the $23$ nondegenerate outer equilibria.
Specifically, for any degree-$6$ harmonic polynomial
$P(X) = H_6(R X) + \sum_{\ell=1}^5 \tilde{H}_\ell(X)$ (where $R \in SO(3)$ is a
free spatial rotation chosen to align $H_6$ favorably with $\tilde{H}_5$), the
rescaled potential $\delta^{-6}(V(\delta X) - V(0))$ converges in
$C^2_{\mathrm{loc}}(\R^3)$ to $c_6 P(X)$ as $\delta \to 0^+$.

To maximize the number of real nondegenerate critical points of $P$ while
keeping its top term fixed to $H_6^{(0)}(R X)$, we build $P$ by \emph{recursive
multi-scale harmonic nesting}.

\subsection{Scale 1: The $100\%$-B\'ezout Quartic Core $P_{27}$}

Consider the $C_{3v}$-symmetric degree-$4$ harmonic polynomial
\begin{equation}\label{eq:p27}
  P_{27}(x,y,z) = Y_{4,0} + 25\,Y_{4,3c} + 15\,Y_{3,3c} - 4\,Y_{3,0} - \tfrac{4}{5}\,Y_{2,0},
\end{equation}
where $Y_{2,0} = 2z^2 - \rho^2$, $Y_{3,0} = 2z^3 - 3\rho^2 z$,
$Y_{3,3c} = x^3 - 3xy^2$,
$Y_{4,0} = 8z^4 - 24\rho^2 z^2 + 3\rho^4$, and
$Y_{4,3c} = z(x^3 - 3xy^2)$.

\begin{proposition}[$100\%$ B\'ezout Reality in Degree $4$]\label{prop:p27}
The gradient system $\nabla P_{27}(x,y,z) = 0$ has B\'ezout bound $3^3 = 27$,
and all $27$ solutions are real, finite, and nondegenerate, with Morse counts
$c_1 = 16$, $c_2 = 11$ ($c_2 - c_1 = -5$).
\end{proposition}

\begin{proof}[Mechanism and certification]
In cylindrical coordinates $(\rho,\theta,z)$, the azimuthal derivative factors
as $\partial_\theta P_{27} = -3\rho^3 \sin(3\theta)\,(25z + 15)$. In addition
to $3$ axial critical points ($\rho = 0$) and $18$ mirror-plane critical points
($\sin 3\theta = 0$, six on each of the three mirror planes), the horizontal
plane $z_* = -15/25 = -3/5$ annihilates $\partial_\theta P_{27}$ identically
and yields a positive radius $\rho_* > 0$ with $|\cos 3\theta| < 1$, producing
$6$ off-mirror critical points: $3 + 18 + 6 = 27$. Complete certification
(\texttt{certified = 27/27}, \texttt{complete = true}, all $27$ real) is
executed in \texttt{src/limiting\_poly.rs} via \texttt{pandsysroots~0.34.0}.
\end{proof}

\subsection{Projective Outer-Scale Collinearity Lemma}

Suppose $P_{d-1} = H_{d-1} + \dots + H_1$ is a harmonic polynomial of degree
$d - 1$ with $N_{d-1}$ real nondegenerate critical points, and consider the
nested perturbation
\[
  P_d(X) = P_{d-1}(X) + \varepsilon_d H_d(X), \qquad 0 < \varepsilon_d \ll 1.
\]
All $N_{d-1}$ critical points of $P_{d-1}$ persist at scale $\|X\| = O(1)$.
At the outer scale $X = Y / \varepsilon_d$, the rescaled equation
$\varepsilon_d^{d-1} \nabla P_d(Y/\varepsilon_d) = 0$ converges to
$\nabla H_d(Y) + \nabla H_{d-1}(Y) = 0$, whose nonzero solutions $Y = r\hat{u}$
($\hat{u} \in S^2$, $r \in \R \setminus \{0\}$) satisfy
\begin{equation}\label{eq:collinear}
  r\,\nabla H_d(\hat{u}) + \nabla H_{d-1}(\hat{u}) = 0
  \iff \nabla H_d(\hat{u}) \times \nabla H_{d-1}(\hat{u}) = 0.
\end{equation}
Because $\nabla H_d$ and $\nabla H_{d-1}$ have opposite parity on $S^2$,
replacing $\hat{u} \mapsto -\hat{u}$ sends $r \mapsto -r$ and leaves
$Y = r\hat{u}$ invariant. Hence the outer equilibria are in exact bijection
with the real projective solutions $[\hat{u}] \in \RP^2$ of the $2$-dimensional
collinearity system~\eqref{eq:collinear} (with $\nabla H_d(\hat{u}) \ne 0$,
$\nabla H_{d-1}(\hat{u}) \ne 0$, and $r = -(\nabla H_{d-1}\cdot \nabla H_d)/\|\nabla H_d\|^2$),
whose B\'ezout budget is $(d-1)^3 - (d-2)^3$ ($37$ for $d = 5$, $61$ for
$d = 6$).

\subsection{Scales 2 and 3: Certified Nested Sextic $P_{71}$}

Applying~\eqref{eq:collinear} twice:
\begin{enumerate}
  \item At degree $5$, choosing
  $H_5 = \cos(\theta_5)\,\frac{\sqrt{11}}{8} Y_{5,0} + \sin(\theta_5)\sqrt{\frac{231}{128}}\,Y_{5,3c}$
  with $\theta_5 = 2.0158552862403435$ and $\varepsilon_5 = 0.05$ yields $25$
  real outer solutions of $\nabla H_5 \times \nabla H_4 = 0$ on $\RP^2$,
  producing a degree-$5$ harmonic polynomial $P_{52} = P_{27} + \varepsilon_5 H_5$
  with $27 + 25 = 52$ real critical points ($c_1 = 26$, $c_2 = 26$; note that
  $c_2 - c_1 = \deg(\nabla H_5/\|\nabla H_5\|) \equiv 0$ identically for every
  odd degree since $\nabla H_5$ is even and factors through $\RP^2$, where
  $H_2(\RP^2; \Z) = 0$).
  \item At degree $6$, scanning $R \in SO(3)$ against the exact sextic harmonic
  $H_6^{(0)}$ of Proposition~\ref{prop:design} finds the Euler $ZYZ$ rotation
  $R = R_Z(0) R_Y(\pi/16) R_Z(0)$ for which
  $\nabla(H_6^{(0)} \circ R) \times \nabla H_5 = 0$ admits $19$ real projective
  solutions on $\RP^2$.
\end{enumerate}

\begin{proposition}[Complete \texttt{pandsysroots~0.34.0} Certification of $P_{71}$]\label{prop:p71}
For $\varepsilon_6 = 3 \times 10^{-7}$ (and likewise for $\varepsilon_6 = 3 \times 10^{-8}$),
the degree-$6$ harmonic polynomial
\[
  P_{71}(X) = P_{27}(X) + \varepsilon_5 H_5(X) + \varepsilon_6 H_6^{(0)}(R X)
\]
satisfies $\Delta P_{71} \equiv 0$, and \texttt{pandsysroots~0.34.0} certifies
all $125 = 5^3$ complex B\'ezout critical points (\texttt{complete = true}) in
$62.7\,\mathrm{ms}$, consisting of exactly $71$ real nondegenerate critical
points ($c_1^{\mathrm{in}} = 35$, $c_2^{\mathrm{in}} = 36$,
$c_2^{\mathrm{in}} - c_1^{\mathrm{in}} = +1$) and $27$ complex-conjugate pairs.
\end{proposition}

\section{Pandrosion V31 Radical Enclosures and Exact 512-Bit Certification}\label{sec:cert}

To prove Theorem~\ref{thm:main} unconditionally on a \emph{concrete, finite}
charge configuration without invoking asymptotic limit arguments, we solve the
$35$-equation perturbed moment system at finite scale $\delta = 10^{-13}$ (with
$\varepsilon_6 = 3\times 10^{-8}$) in $180$-digit arithmetic and round all $44$
coordinates and positive weights $(a_{j,1}, a_{j,2}, a_{j,3}, q_j)$ to exact
rationals in $\Z / 10^{90}$. Table~\ref{tab:charges} summarizes the resulting
$11$ strictly positive charges (exact $90$-digit integers are stored verbatim in
\texttt{data/witness\_n11\_94.json}).

\begin{table}[ht]
\centering
\caption{The $11$ strictly positive point charges of Theorem~\ref{thm:main},
ordered by radius $\|a_j\|$ (truncated to $6$ digits for display; exact
rationals in $\Z / 10^{90}$ are embedded in the Rust certifier).}\label{tab:charges}
\small
\begin{tabular}{@{}rccccc@{}}
\toprule
Shell / Rank & $\|a_j\|$ & $q_j > 0$ & $a_{j,1}$ & $a_{j,2}$ & $a_{j,3}$ \\
\midrule
Inner 1  & $0.033859$ & $4.668205\times 10^{-8}$ & $-0.006606$ & $+0.000000$ & $+0.033209$ \\
Mid 2    & $0.132092$ & $1.491172\times 10^{-4}$ & $-0.057399$ & $-0.109148$ & $-0.047333$ \\
Mid 3    & $0.183119$ & $8.439569\times 10^{-4}$ & $-0.078408$ & $+0.149154$ & $-0.071679$ \\
Mid 4    & $0.183119$ & $8.439601\times 10^{-4}$ & $+0.181751$ & $-0.010101$ & $-0.019930$ \\
Outer 5  & $0.601050$ & $3.309598\times 10^{-1}$ & $-0.308293$ & $-0.396248$ & $+0.330461$ \\
Outer 6  & $0.842740$ & $9.891926\times 10^{-1}$ & $-0.494093$ & $+0.442796$ & $+0.519629$ \\
Outer 7  & $0.842740$ & $9.891941\times 10^{-1}$ & $+0.441768$ & $-0.130084$ & $+0.705783$ \\
Outer 8  & $0.924300$ & $1.341980\times 10^{+0}$ & $+0.363572$ & $-0.838471$ & $-0.138252$ \\
Outer 9  & $0.924302$ & $1.341989\times 10^{+0}$ & $-0.837511$ & $-0.103238$ & $-0.377162$ \\
Outer 10 & $1.334406$ & $2.572726\times 10^{+0}$ & $+0.214173$ & $-0.077375$ & $-1.314831$ \\
Outer 11 & $1.555080$ & $5.069027\times 10^{+0}$ & $+0.830485$ & $+1.305148$ & $-0.158605$ \\
\bottomrule
\end{tabular}
\end{table}

\subsection{Pandrosion V31 Order-$4$ Two-Sided Radical Enclosures}

Evaluating $\nabla V$ and $\Hess V$ requires rigorous enclosures of the $11$
Euclidean distances $r_j = \sqrt{\|x - a_j\|^2}$ over arbitrary-precision
interval boxes without relying on hardware floating-point or \texttt{libm}
square roots. From the Pandrosion~V31 projective support~\cite{PandrosionV31}
specialized to degree $p = 2$, we obtain an unconditionally exact two-sided
rational enclosure.

\begin{lemma}[Pandrosion V31 Projective Two-Sided Square-Root Enclosure]\label{lem:pandrosion}
For every $S > 0$ and every positive approximation $y > 0$,
\begin{equation}\label{eq:pandrosion-sqrt}
  \underline{\Phi}_{\mathrm{proj}}(S, y)
  := \frac{4 y S (y^2 + S)}{S^2 + 6 y^2 S + y^4}
  \;\le\; \sqrt{S} \;\le\;
  \frac{S^2 + 6 y^2 S + y^4}{4 y (y^2 + S)}
  =: \overline{\Phi}_{\mathrm{proj}}(S, y).
\end{equation}
\end{lemma}

\begin{proof}
Direct algebraic factorization gives
\[
  \overline{\Phi}_{\mathrm{proj}}(S, y) - \sqrt{S}
  = \frac{(\sqrt{S} - y)^4}{4 y (y^2 + S)} \ge 0,
  \qquad
  \sqrt{S} - \underline{\Phi}_{\mathrm{proj}}(S, y)
  = \frac{\sqrt{S}\,(\sqrt{S} - y)^4}{S^2 + 6 y^2 S + y^4} \ge 0,
\]
with quartic local convergence $O(|y - \sqrt{S}|^4)$.
\end{proof}

In \texttt{src/dyadic.rs}, every interval $[S_{\mathrm{lo}}, S_{\mathrm{hi}}]$
with $S_{\mathrm{lo}} > 0$ is enclosed at $512$-bit precision by refining a
pure-integer seed through four projective steps, evaluating
$\underline{\Phi}_{\mathrm{proj}}(S_{\mathrm{lo}}, y_{\mathrm{lo}})$ with
numerator rounded strictly down and denominator strictly up, evaluating
$\overline{\Phi}_{\mathrm{proj}}(S_{\mathrm{hi}}, y_{\mathrm{hi}})$ with
numerator rounded strictly up and denominator strictly down, and independently
asserting $r_{\mathrm{lo}}^2 \le S_{\mathrm{lo}}$ and
$r_{\mathrm{hi}}^2 \ge S_{\mathrm{hi}}$ in exact arbitrary-precision integer
arithmetic.

\subsection{3D Interval Krawczyk Certification}

For each candidate center $\hat{x}_k \in \R^3$ ($k = 1,\dots,94$), our Rust
verifier constructs a cube $B_k = \hat{x}_k + [-\rho_k, \rho_k]^3$, computes a
$512$-bit midpoint preconditioner $Y_k \approx (\Hess V(\hat{x}_k))^{-1}$,
evaluates outward-rounded interval enclosures $\nabla V(\hat{x}_k)$ and
$\Hess V(B_k)$ via Lemma~\ref{lem:pandrosion}, and verifies:
\begin{enumerate}
  \item \emph{Pole separation:} $S_{k,j}^{\mathrm{lo}} > 0$ for all $j = 1,\dots,11$;
  \item \emph{Contraction:} $\|M_k\|_\infty < 1$, where $M_k = I - Y_k \Hess V(B_k)$;
  \item \emph{Strict Krawczyk inclusion:}
  \[
    K(B_k) := \hat{x}_k - Y_k \nabla V(\hat{x}_k) + M_k(B_k - \hat{x}_k)
    \;\subset\; \operatorname{int}(B_k);
  \]
  \item \emph{Exact nondegeneracy and Morse index:} $0 \notin \det \Hess V(B_k)$,
  certifying index $1$ when $\det \Hess V(B_k) < 0$ and index $2$ when
  $\det \Hess V(B_k) > 0$;
  \item \emph{Global distinctness:} all $94$ certified boxes $B_1,\dots,B_{94}$
  are pairwise disjoint, and all $11$ exact charge enclosures are pairwise
  disjoint with $q_j^{\mathrm{lo}} > 0$.
\end{enumerate}

By the Krawczyk--Moore--Miranda fixed-point theorem, conditions (1)--(5)
rigorously prove the existence of $94$ distinct nondegenerate critical points
of $V$ in $\R^3 \setminus \{a_1,\dots,a_{11}\}$.

\section{Certified Results and Four-Scale Breakdown}\label{sec:results}

Executing \texttt{cargo run --release} and \texttt{cargo test --release} in the
companion crate \texttt{maxwell-n11-94} certifies all three assertions
($P_{27}$, $P_{71}$, and the full $94$-point 3D Coulomb configuration) in under
$0.11\,\mathrm{s}$ total wall-clock time on standard hardware:
\begin{itemize}
  \item \textbf{Core $P_{27}$ (\texttt{pandsysroots~0.34.0}):} $27/27$ certified,
  \texttt{complete = true}, all $27$ real ($c_1 = 16$, $c_2 = 11$, $c_2 - c_1 = -5$).
  \item \textbf{Limiting $P_{71}$ (\texttt{pandsysroots~0.34.0}):} $125/125$
  certified, \texttt{complete = true}, $71$ real ($c_1 = 35$, $c_2 = 36$,
  $c_2 - c_1 = +1$), verified in $62.70\,\mathrm{ms}$.
  \item \textbf{Direct 3D Coulomb Potential ($512$-bit Pandrosion~V31 + Krawczyk):}
  all $94$ critical points certified in $41.50\,\mathrm{ms}$ with
  $\max_k \|I - Y_k \Hess V(B_k)\|_\infty = 1.386 \times 10^{-21} \ll 1$,
  minimum pole distance $0.033859$, maximum inner box radius
  $2.061 \times 10^{-67}$, maximum outer box radius $5.398 \times 10^{-79}$,
  all $94$ boxes pairwise disjoint, and exact Morse counts $c_1 = 42$,
  $c_2 = 52$ ($c_2 - c_1 = +10 = 11 - 1$).
\end{itemize}

\begin{table}[ht]
\centering
\caption{Four-scale hierarchy of the $94$ certified nondegenerate critical
points of Theorem~\ref{thm:main} and their exact cumulative Poincar\'e--Hopf
indices.}\label{tab:scales}
\footnotesize
\begin{tabular}{@{}clrrrc@{}}
\toprule
Scale & Structural Origin & Count & $(c_1, c_2)$ & $c_2 - c_1$ & Radial Range $\|x\|$ \\
\midrule
1 & $100\%$-B\'ezout $C_{3v}$ quartic core $P_{27}$ & $27$ & $(16, 11)$ & $-5$ & $[7.20\times 10^{-46},\; 1.85\times 10^{-13}]$ \\
2 & Quintic nesting ($P_{27} + \varepsilon_5 H_5$) & $25$ & $(10, 15)$ & $+5$ & $[1.92\times 10^{-12},\; 1.56\times 10^{-11}]$ \\
3 & Sextic nesting vs.\ $11$-charge $H_6^{(0)}$ & $19$ & $(9, 10)$ & $+1$ & $[8.33\times 10^{-9},\; 5.65\times 10^{-8}]$ \\
4 & Multi-shell $n=11$ Kelvin $5$-design outer & $23$ & $(7, 16)$ & $+9$ & $[3.60\times 10^{-2},\; 8.01\times 10^{-1}]$ \\
\midrule
\multicolumn{2}{@{}l}{\textbf{Total (exact 11-charge potential $V$)}} & \textbf{94} & \textbf{(42, 52)} & \textbf{+10} & \textbf{Min pole dist $= 0.033859$} \\
\bottomrule
\end{tabular}
\end{table}

\begin{remark}[Higher Recursive Nesting]
Without constraining $H_6$ to a fixed $11$-charge Kelvin $5$-design, the same
$C_{3v}$ recursive nesting mechanism immediately produces an $89/125$ real
degree-$6$ harmonic polynomial $P_{89}$ ($c_1 = 48$, $c_2 = 41$,
$c_2 - c_1 = -7$) and a $120/216$ real degree-$7$ harmonic polynomial
$P_{120}$, both certified complete by \texttt{pandsysroots~0.34.0}. Realizing
negative-index sextics such as $P_{89}$ or septics such as $P_{120}$ via higher
multi-shell positive Kelvin designs is a promising route toward further
improvements for $n \ge 12$.
\end{remark}

\begin{thebibliography}{9}

\bibitem{ABK2026}
P.~Arathoon, G.~Ball, and M.~D.~Kvalheim,
\emph{The Maxwell Conjecture is False},
preprint, arXiv:2607.27197, July 2026.

\bibitem{PandrosionV31}
I.~Besevic,
\emph{Geometric Root Calculators: Radical supports of orders five and nine and
real-degree enclosures in the Pandrosion rectangle},
preprint V31, November 2026, \url{https://github.com/ivan-fr/pandrosion}.

\bibitem{pandsysroots}
I.~Besevic,
\emph{\texttt{pandsysroots} 0.34.0: Certified polynomial system solver via
simultaneous Pandrosion rectangle sweeps and residue bounds},
crates.io, 2026, \url{https://crates.io/crates/pandsysroots}.

\bibitem{GNS2007}
A.~Gabrielov, D.~Novikov, and B.~Shapiro,
\emph{Mystery of point charges},
Proc.\ Lond.\ Math.\ Soc.\ \textbf{95} (2007), 443--472.

\bibitem{Kintali2026pos}
S.~Kintali,
\emph{Positive Counterexamples to the Maxwell Bound for Four, Six, Eight, Nine,
and Ten Charges},
preprint, August 2026.

\bibitem{Kintali2026quad}
S.~Kintali,
\emph{A Quadratic Lower Bound for Maxwell's Equilibrium Problem},
preprint, August 2026.

\bibitem{Kintali2026sup}
S.~Kintali,
\emph{A Superlinear Lower Bound for Maxwell's Equilibrium Problem},
preprint, August 2026.

\end{thebibliography}

\end{document}
